\subsection{Keep the event's intent when a method is specialized}
\label{sec:authored-event-contract}

A consistent closing book still cannot show whether a collection was permitted
at the time it occurred. In the constructed lifecycle, an assessment on
31 January closes a window and advances its cursor to 1 February. A collection
dated 31 January can satisfy ordinary nondecreasing book dates and all claim
state invariants, yet violate that sealed-window policy. The retained lifecycle
refuses it. A separately selected event contract is needed to preserve that
meaning when the method is authored elsewhere.

Customer and interval assertions expose another boundary. The original full
event states which claim and customer a receipt concerns, or which start and
end dates an assessment covers. A specialized method may embed those values,
so they no longer appear among its declared calculation inputs. Removing them
from the projected event is acceptable only if the original intent remains
available wherever eligibility and specialization must be checked. A method
that happens to produce a consistent state cannot supply the missing assertion
about which customer the caller meant.

\code{ContractedAuthoredCompanyProgram/1} therefore retains every complete
original event beside the projected event. Its caller explicitly selects the
retained receivable lifecycle contract, identified by the native implementation
hash, roles, unit-cost pointer and policy. Before an event, that named owner
checks the full input population, unknowns and applicability on the actual
current book, then derives the intended method under its existing claim,
customer and interval rules. The complete authored materialized request must
equal that intended request after the declared input projection and permitted
scalar-collection transition conversion. Both requests and the removed input
values remain inspectable.

This contract is deliberately specific. The half-payment method previously
authored is a valid alternative algorithm within its stated scope. It does not
satisfy a caller's selection of the original full-payment lifecycle, so that
selection refuses it. An explicit policy choice resolves the difference; a
balanced journal or a valid type cannot choose between the two algorithms.
The native specialization owner is named, not hidden behind an arbitrary
callback or presented as a general declarative financial compiler.

Both eleven-command examples preserve complete states, statements, period
carry and archives, and every full-intent native event request matches the
retained lifecycle. Eighteen refusal comparisons across the two charts cover
customer mismatch, missing or repeated claims, unknown inputs, false
applicability, inconsistent assessment intervals, an early due date and the
sealed-window event. Earlier successful commands remain unchanged when a later
event refuses. Changed projected values and a method inconsistent with the
selected intent are also refused; forged intent evidence fails recomputation.

These comparisons extend the checked refusal scope without proving every
lifecycle rule or supplying actual company facts. General portable authoring
of event contracts, wider financial operations and source/accounting admission
remain unresolved.

The source boundary and counterexamples are retained under
\path{data/accounting-master/authored-event-prerequisite-audit-20261026/}.
The selected-contract evidence is retained under
\path{data/accounting-master/authored-event-contract-20261026/}.

\par
